\chapter{Mạnh Tử — Tính Thiện, Lòng Trắc Ẩn, Dân Vi Quý}
\markboth{Mạnh Tử — Tính Thiện, Lòng Trắc Ẩn, Dân Vi Quý}{Mencius — Original Goodness, Compassion, People First}

\section{Tính Người Vốn Thiện Như Nước Chảy Xuống}
\begin{truyen}{Tính Người Vốn Thiện Như Nước Chảy Xuống}{Human Nature Is Originally Good Like Water Flowing Downward}
\chuhoa{C}{áo Tử tranh luận với Mạnh Tử: ``Tính người không thiện không ác, như nước — đổ về đông thì chảy đông, đổ về tây thì chảy tây.''}
\textit{(Gaozi argued with Mencius: ``Human nature is neither good nor bad, like water — direct it east and it flows east, direct it west and it flows west.'')}

Mạnh Tử đáp: ``Nước thì không phân đông tây, nhưng có phân trên dưới. Tính thiện của người cũng như nước chảy xuống — không ai không thiện, như không có nước nào không chảy xuống.''
\textit{(Mencius replied: ``Water does not distinguish east from west, but it does distinguish up from down. Human goodness is like water flowing downward — there is no person who is not good, just as there is no water that does not flow downward.'')}

``Nhưng có thể đập nước bắn lên trán người — đó có phải tính của nước không? Bắt nước chảy ngược lên núi được — nhưng đó không phải tính của nước. Người làm ác cũng vậy — là vì hoàn cảnh cưỡng bức, không phải vì tính gốc.''
\textit{(``But you can splash water up to a man's forehead — is that the nature of water? You can make water flow up a mountain — but that is not the nature of water. When people do evil, it is likewise because of circumstances forcing it, not because of their original nature.'')}

Mẹ Mạnh Tử chuyển nhà ba lần để chọn môi trường tốt cho con — bởi vì bà hiểu: tính thiện cần môi trường để phát triển.
\textit{(Mencius's mother moved their home three times to choose a good environment for her son — because she understood: goodness of nature needs the right environment to develop.)}

\end{truyen}

\begin{baihoc}
Con người sinh ra với mầm thiện. Nhưng mầm thiện cần được vun xới bởi giáo dục và môi trường đúng. Đây là lý do giáo dục không bao giờ là xa xỉ phẩm.
\textit{(People are born with seeds of goodness. But those seeds need to be nurtured by proper education and environment. This is why education is never a luxury.)}
\end{baihoc}
\ngancach

\section{Tứ Đoan — Bốn Mầm Đức Hạnh}
\begin{truyen}{Tứ Đoan — Bốn Mầm Đức Hạnh}{The Four Sprouts of Virtue}
\chuhoa{M}{ạnh Tử dạy: người ta ai cũng có bốn mầm đức hạnh bẩm sinh, như bốn chi trên thân thể.}
\textit{(Mencius taught: everyone has four innate sprouts of virtue, like the four limbs of the body.)}

``Thấy đứa trẻ sắp ngã xuống giếng, ai cũng có lòng kinh hãi thương xót — đó là mầm của Nhân.''
\textit{(``Seeing a child about to fall into a well, everyone feels alarm and compassion — that is the sprout of Benevolence.'')}

``Biết xấu hổ về điều xấu mình làm — đó là mầm của Nghĩa.''
\textit{(``Feeling shame about the wrong things one does — that is the sprout of Righteousness.'')}

``Biết nhường nhịn — đó là mầm của Lễ.''
\textit{(``Knowing to defer and yield — that is the sprout of Ritual Propriety.'')}

``Phân biệt đúng sai — đó là mầm của Trí.''
\textit{(``Distinguishing right from wrong — that is the sprout of Wisdom.'')}

``Ai có bốn mầm này mà nói mình không thể thiện là tự hại mình. Ai nói vua mình không thể thiện là tự hại vua mình.''
\textit{(``Whoever has these four sprouts and says they cannot be good is injuring themselves. Whoever says their ruler cannot be good is injuring their ruler.'')}

\end{truyen}

\begin{baihoc}
Bạn không cần học đạo đức từ đầu — bạn chỉ cần nhận ra và nuôi dưỡng những mầm đức hạnh đã sẵn có trong mình từ khi sinh ra.
\textit{(You do not need to learn ethics from scratch — you only need to recognize and nurture the sprouts of virtue already present in you from birth.)}
\end{baihoc}
\ngancach

\section{Dân Vi Quý, Xã Tắc Thứ Chi, Quân Vi Khinh}
\begin{truyen}{Dân Vi Quý, Xã Tắc Thứ Chi, Quân Vi Khinh}{The People Are Most Important; the State Comes Next; the Ruler Least}
\chuhoa{Đ}{ây là câu nói táo bạo nhất của Mạnh Tử — và nguy hiểm nhất với các bạo vương thời ông.}
\textit{(This is Mencius's boldest statement — and the most dangerous for the tyrants of his time.)}

``Dân vi quý, xã tắc thứ chi, quân vi khinh'' — Dân quan trọng nhất, đất nước đứng thứ, vua là nhẹ nhất.
\textit{(``The people are most important; the altars of land and grain come next; the ruler is least important.'')}

Một vua hỏi: ``Vậy vua có quyền gì?''
\textit{(A king asked: ``Then what authority does the ruler have?'')}

Mạnh Tử đáp: ``Thiên tử mất lòng dân thì mất thiên hạ. Chư hầu mất lòng thiên tử thì mất nước. Đại phu mất lòng chư hầu thì mất ấp. Dân là nền tảng của quốc gia — nền bền thì nước yên.''
\textit{(Mencius replied: ``When the Son of Heaven loses the favor of the people, he loses the realm. When the feudal lords lose the favor of the Son of Heaven, they lose their states. When the ministers lose the favor of the feudal lords, they lose their domains. The people are the foundation of the nation — when the foundation is solid, the nation is at peace.'')}

Tư tưởng này xuất hiện hơn 2000 năm trước khi các triết gia châu Âu phát triển khái niệm chủ quyền nhân dân.
\textit{(This idea appeared more than 2000 years before European philosophers developed the concept of popular sovereignty.)}

\end{truyen}

\begin{baihoc}
Quyền lực chỉ hợp pháp khi phục vụ người dân. Mọi tổ chức — nhà nước, công ty hay gia đình — vận hành tốt nhất khi đặt lợi ích của những người phụ thuộc vào nó lên trước.
\textit{(Power is only legitimate when it serves the people. Every organization — state, company, or family — functions best when it places the interests of those who depend on it first.)}
\end{baihoc}
\ngancach

\section{Mẹ Mạnh Tử Chuyển Nhà Ba Lần}
\begin{truyen}{Mẹ Mạnh Tử Chuyển Nhà Ba Lần}{Mencius's Mother Moved Three Times}
\chuhoa{T}{huở nhỏ, Mạnh Tử và mẹ sống gần nghĩa địa. Mạnh Tử thường chơi trò giả vờ chôn cất người chết.}
\textit{(In his youth, Mencius and his mother lived near a cemetery. Mencius would play at imitating burial rites.)}

Mẹ ông lo lắng và chuyển nhà đến gần chợ. Mạnh Tử bắt đầu chơi trò bán hàng và mặc cả.
\textit{(His mother was concerned and moved near a market. Mencius began playing at buying and selling.)}

Bà chuyển nhà lần thứ ba, đến gần trường học. Mạnh Tử bắt đầu chơi trò học lễ nghi và tranh luận.
\textit{(She moved a third time, to near a school. Mencius began playing at learning rites and debating.)}

``Đây là chỗ xứng đáng cho con ta,'' bà nói và ở lại.
\textit{(``This is a proper place for my son,'' she said and stayed.)}

Câu chuyện này là nền tảng của triết lý giáo dục Mạnh Tử: môi trường định hình tính cách. Đứa trẻ không chọn được môi trường — người lớn có trách nhiệm chọn thay.
\textit{(This story is the foundation of Mencius's educational philosophy: environment shapes character. Children cannot choose their environment — adults bear the responsibility to choose for them.)}

\end{truyen}

\begin{baihoc}
Chúng ta đều là sản phẩm một phần của môi trường mình lớn lên. Hãy có ý thức về môi trường bạn tạo ra cho con trẻ — hoặc cho chính bạn.
\textit{(We are all partly products of the environment we grew up in. Be intentional about the environment you create for children — or for yourself.)}
\end{baihoc}
\ngancach

\section{Kẻ Cùng Đường Không Thụ Diều Rẻo}
\begin{truyen}{Kẻ Cùng Đường Không Thụ Diều Rẻo}{The Proud Poor Man}
\chuhoa{M}{ạnh Tử kể câu chuyện một người đàn ông nghèo, vợ và thiếp đều kính trọng, tưởng ông là người đáng trọng.}
\textit{(Mencius told the story of a poor man whose wife and concubine both respected him, thinking he was worthy of regard.)}

Nhưng bí mật là: mỗi ngày ông lẻn ra nghĩa địa xin đồ ăn thừa từ người đến cúng tế, rồi về nhà phô trương như đã ăn với người giàu có.
\textit{(But the secret was: each day he would sneak to the cemetery to beg leftover food from those who came to make offerings, then return home bragging about dining with the wealthy.)}

Vợ và thiếp phát hiện ra sự thật, hai người ôm nhau khóc.
\textit{(When his wife and concubine discovered the truth, they wept together.)}

Mạnh Tử kết luận: ``Nếu quan sát cách người đời tìm kiếm giàu sang danh vọng, ít người không làm khiến vợ thiếp phải xấu hổ và khóc.''
\textit{(Mencius concluded: ``If we observe how people in the world seek wealth and honor, there are few whose wives and concubines would not be ashamed and weep.'')}

Câu chuyện không phê phán người nghèo — nó phê phán sự giả tạo và cái giá của danh dự lừa dối.
\textit{(The story does not criticize poverty — it criticizes pretense and the price of false honor.)}

\end{truyen}

\begin{baihoc}
Danh dự xây trên sự lừa dối không phải danh dự — đó là gánh nặng. Sự thật, dù khiêm nhường, nhẹ hơn mọi lời nói dối hào nhoáng.
\textit{(Honor built on deception is not honor — it is a burden. Truth, however humble, is lighter than any glamorous lie.)}
\end{baihoc}
\ngancach

\section{Thiên Tướng Đại Nhiệm Tất Tiên Khổ Kỳ Tâm Chí}
\begin{truyen}{Thiên Tướng Đại Nhiệm Tất Tiên Khổ Kỳ Tâm Chí}{Heaven Prepares Great Tasks Through Suffering}
\chuhoa{M}{ạnh Tử dạy: ``Khi trời muốn giao trọng trách lớn cho người nào, ắt trước làm khổ tâm chí người ấy, làm mỏi mệt gân cốt, làm đói thân thể, làm thiếu thốn thân người, làm rối loạn mọi việc người làm.''}
\textit{(Mencius taught: ``When heaven is about to confer a great responsibility on someone, it first afflicts their mind with suffering, tires their sinews and bones, starves their body, subjects them to poverty, and confounds their undertakings.'')}

``Tất cả để thúc đẩy tâm họ kiên nhẫn, và tăng thêm năng lực mà trước đây họ không có.''
\textit{(``All of this is to stimulate their mind to endurance, and to develop capabilities they did not previously possess.'')}

Ông liệt kê những thánh nhân vĩ đại đều qua gian khổ trước khi trở thành vĩ đại: Shun từ ruộng đồng, Phó Duyệt từ công trường xây dựng, Giao Cách từ cá muối, Quản Trọng từ lao ngục.
\textit{(He listed the great sages who all passed through hardship before becoming great: Shun from the fields, Fu Yue from construction sites, Jiao Ge from salt fish selling, Guan Zhong from prison.)}

Đây không phải triết học của sự đau khổ — mà là quan sát rằng khó khăn cưỡng bức sự tăng trưởng mà thoải mái không thể tạo ra.
\textit{(This is not a philosophy of suffering — but an observation that hardship forces growth that comfort cannot produce.)}

\end{truyen}

\begin{baihoc}
Đừng cầu xin cuộc đời dễ dàng hơn. Hãy cầu mình đủ mạnh hơn. Khó khăn không phải dấu hiệu bạn thất bại — nó có thể là dấu hiệu trời đang chuẩn bị bạn.
\textit{(Do not pray for an easier life. Pray to be stronger. Hardship is not a sign you are failing — it may be a sign that heaven is preparing you.)}
\end{baihoc}
\ngancach

\section{Sinh Ư Ưu Hoạn, Tử Ư An Lạc}
\begin{truyen}{Sinh Ư Ưu Hoạn, Tử Ư An Lạc}{Life Springs from Sorrow; Death Comes from Ease}
\chuhoa{M}{ột học trò hỏi Mạnh Tử: ``Tại sao những người sống khổ lại thường mạnh và sống lâu hơn người sống sung sướng?''}
\textit{(A student asked Mencius: ``Why do people who live hard lives often seem stronger and live longer than those who live comfortably?'')}

Mạnh Tử trả lời bằng quan sát: ``Sinh ở lo âu, chết ở an nhàn.''
\textit{(Mencius responded with an observation: ``Life springs from sorrow and calamity; death comes from ease and pleasure.'')}

``Quốc gia không có kẻ thù bên ngoài và mối lo bên trong thường diệt vong. Cá nhân không có sự kháng cự và thử thách trong cuộc đời thường yếu đi.''
\textit{(``A state that has no outside enemies and internal troubles will usually perish. An individual who has no resistance and challenges in life usually weakens.'')}

Cơ thể tập thể dục trở nên mạnh vì phải vượt qua kháng cự. Trí tuệ phát triển khi đối mặt với vấn đề. Tính cách rèn trong gian khổ.
\textit{(The body strengthened by exercise becomes strong by overcoming resistance. The mind grows when facing problems. Character is forged in hardship.)}

Sự thoải mái hoàn toàn là kẻ thù của sự tăng trưởng.
\textit{(Complete comfort is the enemy of growth.)}

\end{truyen}

\begin{baihoc}
Hãy cảm ơn những thách thức — chúng giữ cho tâm trí và tính cách bạn sắc nét. Cuộc sống không có ma sát là cuộc sống không phát triển.
\textit{(Be grateful for challenges — they keep your mind and character sharp. A life without friction is a life without growth.)}
\end{baihoc}
\ngancach

\section{Mạnh Tử Và Vua Lương Huệ Vương}
\begin{truyen}{Mạnh Tử Và Vua Lương Huệ Vương}{Mencius and King Hui of Liang}
\chuhoa{V}{ua Huệ của nước Lương gặp Mạnh Tử và hỏi: ``Ngài đến từ xa — có điều gì lợi cho nước ta không?''}
\textit{(King Hui of Liang met Mencius and asked: ``You have come from afar — is there anything that will profit my kingdom?'')}

Mạnh Tử đáp thẳng: ``Vua sao lại hỏi về lợi? Chỉ cần có Nhân Nghĩa là đủ.''
\textit{(Mencius replied directly: ``Why must Your Majesty speak of profit? There need only be benevolence and righteousness.'')}

``Nếu vua hỏi điều gì lợi cho nước, đại thần hỏi điều gì lợi cho nhà, sĩ thứ hỏi điều gì lợi cho thân — trên dưới tranh lợi, nước sẽ nguy.''
\textit{(``If the king asks what profits the state, ministers ask what profits their households, scholars and commoners ask what profits themselves — when all levels compete for profit, the state will be in danger.'')}

Mạnh Tử không phủ nhận tầm quan trọng của lợi ích — ông cảnh báo về nguy hiểm khi lợi ích trở thành giá trị duy nhất.
\textit{(Mencius did not deny the importance of benefit — he warned of the danger when profit becomes the only value.)}

Một xã hội chỉ nói về lợi nhuận sẽ xói mòn sự tin tưởng, lòng trung thành và phẩm giá con người.
\textit{(A society that speaks only of profit will erode trust, loyalty, and human dignity.)}

\end{truyen}

\begin{baihoc}
Hỏi ``điều này có lợi gì không?'' là câu hỏi hợp lý. Nhưng khi đó là câu hỏi duy nhất bạn đặt ra, bạn đã bắt đầu mất đi điều quan trọng hơn lợi nhuận.
\textit{(Asking ``what is the profit in this?'' is a reasonable question. But when it is the only question you ask, you have begun losing something more important than profit.)}
\end{baihoc}
\ngancach

\section{Ngư Dữ Hùng Chưởng — Cá Và Chân Gấu}
\begin{truyen}{Ngư Dữ Hùng Chưởng — Cá Và Chân Gấu}{Fish and Bear's Paw — On Choosing Between Two Goods}
\chuhoa{M}{ạnh Tử nói: ``Cá là điều ta muốn. Chân gấu cũng là điều ta muốn. Nếu không thể có cả hai, ta bỏ cá lấy chân gấu.''}
\textit{(Mencius said: ``Fish is something I desire. Bear's paw is also something I desire. If I cannot have both, I will give up fish and take bear's paw.'')}

``Sống là điều ta muốn. Nghĩa cũng là điều ta muốn. Nếu không thể có cả hai, ta bỏ sống giữ nghĩa.''
\textit{(``Life is something I desire. Righteousness is also something I desire. If I cannot have both, I will give up life and hold onto righteousness.'')}

Học trò ngạc nhiên: ``Thưa thầy, thầy coi nghĩa quan trọng hơn mạng sao?''
\textit{(Students were surprised: ``Master, do you consider righteousness more important than life?'')}

``Không phải ta không ham sống — ta ham sống nhiều lắm. Nhưng có thứ ta ham hơn sống, nên không làm bất cứ điều gì để sống còn. Cũng không phải ta không sợ chết — nhưng có thứ ta sợ hơn chết, nên không tránh được nguy hiểm.''
\textit{(``It is not that I do not desire life — I desire it greatly. But there is something I desire more than life, so I will not do anything to preserve it at any cost. It is not that I do not fear death — but there is something I fear more than death, so I do not always avoid danger.'')}

Điều đó, Mạnh Tử nói, là không thể mất — lương tâm và phẩm giá.
\textit{(That thing, Mencius said, must not be lost — conscience and dignity.)}

\end{truyen}

\begin{baihoc}
Mỗi người có một ``chân gấu'' — điều họ sẽ không đánh đổi dù giá nào. Biết điều đó là gì với bạn nghĩa là bạn đã biết mình là ai.
\textit{(Every person has their ``bear's paw'' — something they will not trade at any price. Knowing what that is for you means you know who you are.)}
\end{baihoc}
\ngancach

\section{Bất Nhẫn Nhân Chi Tâm — Lòng Không Thể Chịu Được Khi Người Khác Khổ}
\begin{truyen}{Bất Nhẫn Nhân Chi Tâm — Lòng Không Thể Chịu Được Khi Người Khác Khổ}{The Heart That Cannot Bear to See Others Suffer}
\chuhoa{M}{ạnh Tử mô tả khoảnh khắc mà ông tin rằng chứng minh bản chất thiện của con người.}
\textit{(Mencius described a moment he believed proves the innate goodness of human nature.)}

``Bất kỳ ai thấy đứa trẻ sắp ngã xuống giếng đều sẽ có lòng kinh hãi và thương xót — không phải vì muốn kết thân với cha mẹ đứa trẻ, không phải vì muốn tiếng khen của làng xóm, không phải vì sợ tai tiếng nếu không cứu.''
\textit{(``Anyone who sees a child about to fall into a well will feel alarm and compassion — not because they want to befriend the child's parents, not because they want praise from the village, not because they fear blame if they do not act.'')}

``Cảm xúc đó là tự nhiên, tức thì, không tính toán — đó chính là bằng chứng của Nhân trong tâm người.''
\textit{(``That feeling is natural, immediate, uncalculated — that is the evidence of Benevolence in the human heart.'')}

Thực nghiệm tâm lý học hiện đại xác nhận: mọi người ở khắp nơi trên thế giới đều có phản ứng cảm xúc tức thì với nỗi đau của người khác — bất kể văn hóa hay giáo dục.
\textit{(Modern psychological experiments confirm: people everywhere in the world have an immediate emotional response to others' pain — regardless of culture or education.)}

Mạnh Tử đúng: tính thiện không phải học được — nó được phát hiện.
\textit{(Mencius was right: goodness is not something learned — it is something discovered.)}

\end{truyen}

\begin{baihoc}
Lòng trắc ẩn không phải yếu đuối — đó là bằng chứng bạn còn người. Đừng để xã hội dạy bạn tắt đi thứ thiên nhiên đặt vào bạn.
\textit{(Compassion is not weakness — it is evidence of your humanity. Do not let society teach you to extinguish what nature placed within you.)}
\end{baihoc}
